% Part: first-order-logic
% Chapter: introduction
% Section: syntax

\documentclass[../../../include/open-logic-section]{subfiles}

\begin{document}

\olfileid{fol}{int}{syn}

\olsection{نحو}

نخست باید دقیق کنیم که کدام رشته‌های نمادها !!{sentence}s منطق مرتبهٔ
اول به شمار می‌آیند. این کار را بعداً انجام خواهیم داد؛ فعلاً تنها با
مثال پیش می‌رویم. اجزای سازندهٔ پایه---واژگان---منطق مرتبهٔ اول به دو
بخش تقسیم می‌شوند. بخش نخست نمادهایی است که با آن‌ها دربارهٔ چیزهای
معین سخن می‌گوییم یا چیزهای معین را مشخص می‌کنیم. چیزها را با
!!{constant}s مشخص می‌کنیم و با !!{predicate}s دربارهٔ چیزهای
مشخص‌شده سخن می‌گوییم. برای نمونه، ممکن است از $\Obj a$ به‌عنوان
!!a{constant} برای مشخص‌کردن یک چیز منفرد استفاده کنیم و سپس با
!!{sentence}~$\Atom{\Obj P}{\Obj a}$ چیزی دربارهٔ آن بگوییم. اگر برای
«$\Obj a$» و «$\Obj P$» معنایی در نظر داشته باشید، می‌توانید
$\Atom{\Obj P}{\Obj a}$ را مانند جمله‌ای در زبان طبیعی بخوانید (و
احتمالاً هنگام نخستین آشنایی با منطق صوری چنین کرده‌اید). پس از داشتن
چنین !!{sentence}s ساده‌ای در منطق مرتبهٔ اول، می‌توانیم با بخش دوم
واژگان، یعنی نمادهای منطقی (پیوندگرها و سورها)، جمله‌های پیچیده‌تر بسازیم.
بنابراین، برای نمونه، می‌توانیم عبارت‌هایی مانند
$(\Atom{\Obj P}{\Obj a} \land \Atom{\Obj Q}{\Obj b})$ یا
~$\lexists[\Obj x][\Atom{\Obj P}{\Obj x}]$ تشکیل دهیم.

برای به‌دست‌دادن تعریف‌های دقیق معناشناسی و قواعد دستگاه‌های
!!{derivation} که پژوهش فرامنطقیِ دقیق به آن‌ها نیاز دارد، نخست باید
دقیقاً تعریف کنیم چه چیزی !!a{sentence} منطق مرتبهٔ اول به شمار می‌آید.
فهم ایدهٔ پایه آسان است: برخی !!{sentence}s ساده را فقط از
!!{predicate}s و !!{constant}s، مانند~$\Atom{\Obj P}{\Obj a}$، تشکیل
می‌دهیم؛ سپس با پیوندگرها و سورها جمله‌های پیچیده‌تر را از آن‌ها
می‌سازیم. اما دقیقاً طبق چه قواعدی مجازیم !!{sentence}s پیچیده‌تر بسازیم؟
این قواعد باید مشخص شوند؛ وگرنه «!!{sentence} منطق مرتبهٔ اول» را
به‌اندازهٔ کافی دقیق تعریف نکرده‌ایم. چند مسئله در میان است. نخست باید
رشته‌های درست را !!{sentence} به شمار آوریم. دوم باید این کار را چنان
انجام دهیم که بتوانیم دربارهٔ \emph{همهٔ} !!{sentence}s برهان‌های ریاضی
بیاوریم. سرانجام، باید برخی عملیات ابتدایی روی !!{sentence}s را نیز
دقیق تعریف کنیم، مانند «هر $\Obj x$ در~$!A$ را با~$\Obj b$ جایگزین کن».
مشکل آن است که سورها و !!{variable}s موجود در منطق مرتبهٔ اول، شیوهٔ
انجام این کار را کاملاً بدیهی نمی‌گذارند. برای نمونه، آیا
$\lexists[\Obj x][\Atom{\Obj P}{\Obj a}]$ باید !!a{sentence} به شمار
آید؟ $\lexists[\Obj x][\lexists[\Obj x][\Atom{\Obj P}{\Obj x}]]$ چطور؟
نتیجهٔ «$\Obj x$ را در $(\Atom{\Obj P}{\Obj x} \land \lexists[\Obj
x][\Atom{\Obj P}{\Obj x}])$ با~$\Obj b$ جایگزین کن» باید چه باشد؟

\end{document}
